\clearpage
\subsection*{Пределы точности и физическая интерпретация}
Измеренный ток насыщения растёт от $16{,}73$ до $67{,}2$~мкА.
При сравнимых энергиях электронов это соответствует росту концентрации
почти в четыре раза (\cref{tab:base}). Рост числа носителей заряда
согласуется с увеличением тока разряда. По трём режимам отдельный закон
зависимости концентрации от тока установить трудно.

Для проверки устойчивости обработки полуширину центрального участка
поочерёдно заменяли на $3$ и $6{,}5$~В, границу крайних участков ---
на $15$ и $18$~В. Менялся один параметр за один расчёт. При токах
$1{,}50$, $3{,}00$ и $4{,}99$~мА энергии при таком выборе составили
$3{,}25$--$3{,}57$, $3{,}33$--$3{,}96$ и $2{,}54$--$3{,}86$~эВ
\ifdefined\AIAppendix
соответственно. До подтверждения рукописи автором дополнительно исключали
десять первоначально спорных отсчётов, что меняло оценки на
$0{,}06$--$0{,}09$~эВ. Автор подтвердил все принятые значения;
в итоговый расчёт входят все 83 точки.
\else
соответственно. Исключение отсчётов с неоднозначно прочитанными цифрами
изменило основные оценки примерно на $0{,}06$--$0{,}09$~эВ.
\fi

Если при тех же крайних участках и фиксированном $\Is$ аппроксимировать
центральные точки прямыми с общим наклоном и двумя разными свободными
членами, получаются $3{,}50$, $4{,}31$ и $5{,}20$~эВ. При максимальном
токе поправка достигает $49\,\%$, существенно превышая стандартную ошибку
основного построения. Журнал не позволяет установить причину сдвига ветвей.
Поэтому зависимость температуры от тока по этим данным неоднозначна.

Для всех режимов $T_e$ существенно выше $T_i$.
Электронная поляризационная длина равна $0{,}058$--$0{,}112$~мм.
Даже длина зонда $5{,}2$~мм превышает её примерно в 46--90 раз, что
согласуется с квазинейтральностью объёма положительного столба на
миллиметровых масштабах. Радиус самой трубки в журнале не указан.
В сфере экранирования находится около 30--59 ионов; по принятому в
описании критерию $N_D\gg1$ запас составляет несколько десятков. Это позволяет использовать приближение идеальной плазмы.
Доля ионизованных атомов остаётся меньше $10^{-6}$ во всех трёх режимах.

\section{Заключение}
В диапазоне $0{,}51$--$5{,}01$~мА получена падающая ВАХ разряда.
\ifdefined\AIAppendix
Наибольший по модулю локальный наклон приведён с расширенным бюджетом неопределённости в \cref{tab:r}
\else
Наибольший по модулю локальный наклон равен $(-31{,}4\pm3{,}3)$~кОм
\fi
при понижении тока, на интервале $0{,}51$--$0{,}77$~мА.
Зондовая диагностика дала концентрацию от $1{,}48\cdot10^{10}$ до
$5{,}80\cdot10^{10}$~см$^{-3}$.
При этом частота отсечки возросла с $1{,}09$ до $2{,}16$~ГГц.
Основная оценка энергии электронов составляет $3{,}3$--$3{,}7$~эВ;
наиболее существенное ограничение её точности связано со смещением
зондового тока при переключении полярности.

\section*{Источники}
\begin{enumerate}
\item \textit{3.5.1 (4.16). Изучение плазмы газового разряда в неоне.
Дополнительное описание}. Редакция от 11 августа 2021~г.
С.~1 (установка), с.~2--3 (измерения), с.~3--4 (обработка).
\item Лабораторный журнал работы 3.5.1 (4.16), 17.11.2026.
С.~1 --- ВАХ разряда и первая полярность серии $4{,}99$~мА;
с.~2--3 --- продолжение зондовых измерений. Первичные числовые данные
доступны по адресу \url{https://github.com/duxihw/labs_lpr/tree/main/3.5.1/data/raw}.
\item Никулин М.~Г., Попов П.~В., Нозик А.~А. и др.
\textit{Лабораторный практикум по общей физике. Т.~2. Электричество и магнетизм}.
Под ред. А.~В.~Максимычева и М.~Г.~Никулина. 2-е изд.
М.~: МФТИ, 2019. С.~283--284, рис.~6. ISBN 978-5-7417-0709-8.
\end{enumerate}
